\documentclass[aspectratio=169,11pt]{beamer}

\definecolor{forestgreen}{RGB}{34, 139, 34}

\usetheme{metropolis}
%\usecolortheme{beaver}

%\usepackage[dvipsnames]{xcolor}

\usepackage{amssymb}

\usepackage{proof}



\newcommand{\inl}{\textsf{inl}}
\newcommand{\inr}{\textsf{inr}}

\newcommand{\K}{^\mathrm{k}}

\newcommand{\tto}{\longrightarrow}
\newcommand{\ttoK}{\longrightarrow\K}


\newcommand{\id}{\textsf{id}}
\newcommand{\seq}{\mathbin{\triangleright}}

\newcommand{\idK}{\mathsf{id}^{\tiny \blacktriangleright}}
\newcommand{\seqK}{\mathbin{\blacktriangleright}}

\newcommand{\jd}{\mathsf{jd}}
\newcommand{\prl}{\parallel}


\newcommand{\ema}[1]{\textcolor{blue}{#1}}
\newcommand{\emb}[1]{\textcolor{forestgreen}{#1}}

\newcommand{\ns}[1]{}


\newcommand{\greenb}[1]{{\color{forestgreen}\boxed{\color{black}{#1}}}}
\newcommand{\redb}[1]{{\color{red}\boxed{\color{black}{#1}}}}

\newcommand{\blue}[1]{\textcolor{blue}{#1}}

\newcommand{\gray}[1]{\textcolor{gray}{#1}}



\title{Effects in parallel, with concurrent monads}

\author{Tarmo Uustalu \\
  Reykjvik University / Tallinn University of Technology}

\date{Lambda World '26, Torremolinos, 29--30 Oct.\ 2026}


\begin{document}


\begin{frame}


\maketitle  

\end{frame}    



\begin{frame}

\frametitle{What this is about}

  In functional programming, \ema{monads} are a structure facilitating
  \emb{sequential composition} of effectful functions

\bigskip
  
  Is there a similar thing for \emb{parallel composition}?

  Yes, but underappreciated as such: \ema{applicatives} \gray{[lax
  monoidal functors]}

\bigskip

  But for \emb{both of the two} at once?

  Yes, and you deserve to know: \ema{concurrent monads}
  

\end{frame}


% \begin{frame}

%   \begin{tabular}{c|c}
%     monads & concurrent \\
%            & monads \\[12pt]
%     \multicolumn{2}{c}{an abstraction for} \\
%     \multicolumn{2}{c}{programming with} \\[12pt]
%     effectful operations & effectful operations \\[12pt]
%     \multicolumn{2}{c}{and} \\[12pt]
%     sequential composition & sequential composition\\
%     ~ & parallel composition
%   \end{tabular}  


% \end{frame}

\begin{frame}[fragile]

\frametitle{Monads in Haskell}

{\scriptsize a bit differently than in Prelude}

\scriptsize
\begin{verbatim}
class Functor t => Mnd t where
  return :: x -> t x
  join :: t (t x) -> t x

newtype Reader s x = R (s -> x)

instance Mnd (Reader s) where
  return x = R (\ _ -> x)
  join (R f) = R (\ s -> let R g = f s in g s)  

newtype State s x = S (s -> (s, x))

instance Mnd (State s) where
  return x = S (\ s -> (s, x))
  join (S f) = S (\ s -> let (s', S g) = f s in g s')    
\end{verbatim}

\end{frame}

\begin{frame}

\frametitle{The mission of a monad}
  
  A monad 
  provides

  \begin{itemize}
    
  \item a combinator for sequential composition \\
  of functions with some kind of effect

  \item guarantees that it is sensible:

  \begin{itemize}

  \item with the right type

  \item with the right equational property\footnote{if the monad is a monad indeed} \\
    \gray{(= associativity)}

  \end{itemize}

  \end{itemize}
  
\end{frame}



% \begin{frame}

% \frametitle{Seq comp, with the right type (for exceptions)} 
  
%   ($E$ -- a fixed type, for exceptions, $+$ -- disjoint sum) 

% \[  
%   \infer*[?]{k \seqK \ell = ? : X \tto Z + E}{
%     k : X \tto \redb{Y + E}
%     &
%     \ell : \redb{Y} \tto Z + E
%   }
% \]

% \[  
%   \infer*[?]{\ns{k \seqK \ell = ? :} X \tto \redb{Z + E}}{
%     k : X \tto \greenb{Y + E}
%     &
%     \infer{\greenb{Y + E} \tto \redb{(Z + E) + E}}{
%       \ell : Y \tto Z + E
%     }   
%   }
% \]


% \end{frame}



% \begin{frame}

% \[  
%   \infer{k \seqK \ell \ns{= k \seq (\ell + E) \seq \mu_Z} : X \tto \greenb{Z + E}}{
%     k : X \tto \greenb{Y + E}
%     &
%     \infer{(\ell + E) \seq \mu_Z : \greenb{Y + E} \tto Z + E}{
%        \infer{\ell + E : Y + E \tto (Z + E) + E}{
%          \ell : Y \tto Z + E
%        }
%        &
%        \ns{\mu_Z =} \blue{[\id_{Z+E}, \inr_{Z,E}]} : (Z + E) + E \tto \greenb{Z + E}
%     }   
%   }  
% \]

% \[
%   \idK_X = \blue{\inl_{X,E}} : X \to X + E
% \]  

% \end{frame}



\begin{frame}

\frametitle{Sequential composition, with the right type}
  
$T$ -- a functor, for some kind of effect

Effectful functions $X \to Y$ are the same as pure functions $X \to T Y$

% \[  
%   \infer={k : X \tto T Y}{
%     k : X \ttoK Y
%   }  
% \]
  
\[    
  \infer*[?]{\ns{k \seqK \ell = ? :} X \tto T Z}{
    k : X \tto \redb{T Y} 
    &
    \ell : \redb{Y} \tto T Z
  }
\]

\[    
  \infer*[?]{\ns{k \seqK \ell = ? :} X \tto \redb{T Z}}{
    k : X \tto \greenb{T Y} 
    &
    \infer{T \ell : \greenb{T Y} \tto \redb{T (T Z)}}{
      \ell : Y \tto T Z
    }  
  }
\]

\end{frame}

\frametitle{Sequential computation, with the right type}

\begin{frame}

\[
  \infer{k \seqK \ell := k \seq T \ell \seq \mu_Z : X \tto T Z}{
    k : X \tto \greenb{T Y}
    &
    \infer{T \ell \seq \mu_Z  : \greenb{T Y} \tto \greenb{T Z}}{
       \infer{T \ell : T Y \tto T (T Z)}{
         \ell : Y \tto T Z
       }
       &
       \blue{\mu_Z : T (T Z) \tto T Z} 
    }   
  }  
\]

\[
  \idK_X := \blue{\eta_X : X \tto T X}
\]
  
\end{frame}


\begin{frame}

\frametitle{Sequential computation, with the right equations}


\[
  \infer{(k \seqK \ell) \seqK m = k \seqK (\ell \seqK m)}{
    \infer{\vdots}{
      \blue{\mu_{TX} \seq \mu_X = T \mu_X \seq \mu_X}
    }
  }  
\]  


\[
  \infer{\idK_X \seqK k = k}{
    \infer{\vdots}{
      \blue{\eta_{TX} \seq \mu_X = \id_{TX}}
    }
  } 
\]

\[
  \infer{k = k \seqK \idK_X}{
    \infer{\vdots}{
      \blue{\id_{TX} = T \eta_X \seq \mu_X}
    }
  } 
\]

\end{frame}  


\begin{frame}

\frametitle{A concurrent monad's mission}

 A concurrent monad provides

  \begin{itemize}

  \item combinators for sequential and parallel composition \\
    of functions with some kind of effect

  \item guarantees that they are sensible:

  \begin{itemize}

  \item with the right typing

  \item with the right (in)equational properties
    \footnote{if the concurrent monad is a concurrent monad indeed} \\
  \gray{(= both monotone, both associative, \\
  interrelating via inequational interchange)}

  \end{itemize}
  \end{itemize}

\end{frame}  

\begin{frame}

\frametitle{Parallel composition, with the right typing}

\[
  \infer*[?]{X \times U \tto \redb{T (Y \times V)}}{
    \infer{k \times \ell : X \times U \tto \redb{T Y \times T V}}{
      k : X \tto T Y
      & 
      \ell : U \tto T V
    }  
  }
\]

\[
  \infer{\emb{k \prl \ell} := k \times \ell \seq \phi_{Y,V} : X \times U \tto T (Y \times V)}{
    \infer{k \times \ell : X \times U \tto T Y \times T V}{
      k : X \tto T Y
      & 
      \ell : U \tto T V
    }
    &
    \blue{\phi_{Y,V} : T Y \times T V \tto T (Y \times V)} 
  }
\]

\[
\emb{\jd} := \blue{\phi^0 : 1 \tto T 1}
\]  

\end{frame}
  

\begin{frame}[fragile]

\frametitle{Concurrent monads in Haskell, reader, state}

\scriptsize
\begin{verbatim}
class Functor t => Lmf t where
  nil :: () -> t ()
  merge :: (t x, t y) -> t (x, y)

class (Monad t, Lmf t) => ConcurMnd t where

newtype Reader s x = R (s -> x)

instance Lmf (Reader s) where
  nil () = R (\ _ -> ())
  merge (R f, R g) = R (\ s -> (f s, g s))

instance ConcurMnd (Reader s)

instance MeetSemi s => Lmf (State s) where
  nil () = S (\ _ -> (top, ()))
  merge (S f, S g) = S (\ s -> let ((s0, x), (s1, y)) = (f s, g s) in (meet (s0, s1), (x, y)))    

instance MeetSemi s => ConcurMnd (State s)

\end{verbatim}

\end{frame}


\begin{frame}[fragile]

  \frametitle{A concurrent monad of resumptions for shared state
    and interleaving}

\scriptsize
  
\begin{verbatim}
data Res s x = Ret x | Grab (s -> [(s, Res s x)])

instance Mnd (Res s) where 
  return x = Ret x
  join (Ret r)  = r
  join (Grab k) = Grab (\ s -> [ (s', join r) | (s', r) <- k s ])

instance Lmf (Res s) where
  nil () = Ret ()
  merge (Ret x, Ret y)  = Ret (x, y)
  merge (Ret x, Grab l) =
    Grab (\ s -> [ (s', merge (Ret x, r)  | (s', r) <- l s ])
  merge (Grab k, Ret y) = 
    Grab (\ s -> [ (s', merge (r, Ret y)  | (s', r) <- k s ])
  merge (Grab k, Grab l) =
    Grab (\ s -> [ (s', merge (r, Grab l) | (s', r) <- k s ]
              ++ [ (s', merge (Grab k, r) | (s', r) <- l s ])
\end{verbatim}
  
\end{frame}


\begin{frame}

\frametitle{We saw}

  \ema{Parallel composition} is as fundamental as sequential
  composition, \\ also for effectful functions

  \ema{Monads} axiomatize only sequential composition

\bigskip  

  \ema{Concurrent monads} axiomatize both compositions
  
  The two compositions interrelate via \ema{inequational interchange}

  This presupposes an \ema{order} on functions between the
  same types

\end{frame}


\begin{frame}

\frametitle{We also learned}


  Concurrent monads relax \ema{commutative monads}, \\
  which have equational interchange

\bigskip
  
  Concurrent monads differ from \ema{strong monads}, \\
  which support parallel
  composition only of an effectful function and a pure
  function, \\ with equational interchange

  
\end{frame}
  

\begin{frame}

\frametitle{We did not see}

  Concurrent monads are concurrent monoids in the right place \\
  \gray{(a
  suitable ordered duoidal category)}

  Concurrent monads are (2-categorical) monads in the right place \\
  \gray{(a suitable 2-category)}

  Concurrent monads are meaningfully definable also for orders
  between effectful functions only
  
\end{frame}



\begin{frame}

\frametitle{For more information}

  E Rivas, M Jaskelioff, preprint 2019

  H Paquet, P Saville, LICS '24

  E Rivas, T Uustalu, PPDP '24

  F Breuvart, H Paquet, FSCD '25

  B Gunnarsson, U of Iceland MSc thesis '26

\end{frame}  



\end{document}